% GENERATED FILE. Edit canonical lessons, then run npm run generate:books.
% canonical-source-hash: fnv1a64:b83ea56082157920
% canonical-lessons: ML-C455-ormmippikkuka, ML-C455-bahumanikkuka, ML-C455-tulyamaya, ML-C455-dharmmikamaya, ML-C455-upadravikkuka

\chapter{Respect and Fairness}
\label{ch:ml-respect-and-fairness}
\hlchaptermodality{455}

\begin{chapteropening}
\textbf{By the end of this chapter:} \emph{I can say to remind, to respect, equal, moral and to harm.}

Complete the last lesson of chapter 455: I can say to remind, to respect, equal, moral and to harm.
\end{chapteropening}

\section[\texorpdfstring{\ml{ഓർമ്മിപ്പിക്കുക} (ōrmmippikkuka)}{ഓർമ്മിപ്പിക്കുക (ōrmmippikkuka)}]{\ml{ഓർമ്മിപ്പിക്കുക} (ōrmmippikkuka) — to remind}

\label{lesson:ML-C455-ormmippikkuka}



\begin{quote}
Before the new one: say the Malayalam for to be convinced, then the Malayalam for to confirm.
\end{quote}

\subsection*{You'll want to know: \ml{ഓർമ്മിപ്പിക്കുക}}
\textbf{\ml{ഓർമ്മിപ്പിക്കുക}} — \emph{ōrmmippikkuka} — \textquotedblleft{}to remind\textquotedblright{}.

\begin{quote}
\ml{ഫീസ്} \ml{അടയ്ക്കാൻ} \ml{അമ്മ} \ml{എന്നെ} \ml{ഓർമ്മിപ്പിച്ചു}.
\end{quote}

\begin{quote}
\emph{phīsŭ aṭaykkān amma enne ōrmmippiccu.}
\end{quote}

\begin{quote}
Mother reminded me to pay the fee.
\end{quote}

Past: \ml{ഓർമ്മിപ്പിച്ചു}, \emph{ōrmmippiccu}.

The causative of \ml{ഓർമ്മിക്കുക}, to remember: to make someone remember.

\begin{cousinweb}
\textbf{\ml{ഓർമ്മിപ്പിക്കുക}} (\emph{ōrmmippikkuka}) is native Dravidian, the causative of \ml{ഓർമ്മിക്കുക}, \textquotedblleft{}to remember\textquotedblright{}, on \ml{ഓർമ്മ}, \textquotedblleft{}memory\textquotedblright{} (DEDR 1059). DEDR also lists \emph{ōrpikkuka}, \textquotedblleft{}to remind\textquotedblright{}.
\end{cousinweb}

\begin{grammarlens}[title={What you've built}]
One more word for saying what happens if.
\end{grammarlens}

\subsection*{Your turn}
\begin{itemize}
\raggedright
  \item \emph{Say it:} \emph{ōrmmippikkuka}
  \item \emph{Say it:} \emph{ōrmmippikkuka}, once more
  \item \emph{Say it:} \emph{phīsŭ aṭaykkān amma enne ōrmmippiccu.} — Mother reminded me to pay the fee.
  \item \emph{Recall:} say the Malayalam for to be convinced, then the Malayalam for to confirm, then say \emph{ōrmmippikkuka} again
\end{itemize}

\begin{tcolorbox}[breakable,colback=teal!4,colframe=teal!35!black,title={Before you move on}]
What does \ml{ഓർമ്മിപ്പിക്കുക} mean? (\textquotedblleft{}To remind\textquotedblright{}.) Say it once more.
\end{tcolorbox}
\section[\texorpdfstring{\ml{ബഹുമാനിക്കുക} (bahumānikkuka)}{ബഹുമാനിക്കുക (bahumānikkuka)}]{\ml{ബഹുമാനിക്കുക} (bahumānikkuka) — to look up to; to respect}

\label{lesson:ML-C455-bahumanikkuka}



\begin{quote}
Before the new one: say the Malayalam for to confirm, then the Malayalam for to remind.
\end{quote}

\subsection*{You'll want to know: \ml{ബഹുമാനിക്കുക}}
\textbf{\ml{ബഹുമാനിക്കുക}} — \emph{bahumānikkuka} — \textquotedblleft{}to look up to; to respect\textquotedblright{}.

\begin{quote}
\ml{കുട്ടികൾ} \ml{ആ} \ml{അധ്യാപകനെ} \ml{വളരെ} \ml{ബഹുമാനിക്കുന്നു}.
\end{quote}

\begin{quote}
\emph{kuṭṭikaḷ ā adhyāpakane vaḷare bahumānikkunnu.}
\end{quote}

\begin{quote}
The children respect that teacher very much.
\end{quote}

Past: \ml{ബഹുമാനിച്ചു}, \emph{bahumāniccu}.

The noun is \ml{ബഹുമാനം}, respect.

\begin{cousinweb}
\textbf{\ml{ബഹുമാനിക്കുക}} (\emph{bahumānikkuka}) is built from Sanskrit \emph{bahumāna}, \textquotedblleft{}high esteem\textquotedblright{} (CDIAL 9192), \emph{bahu}, \textquotedblleft{}much\textquotedblright{}, and \emph{māna}, \textquotedblleft{}regard\textquotedblright{}, with the Malayalam verb ending -\ml{ിക്കുക}.
\end{cousinweb}

\begin{grammarlens}[title={What you've built}]
One more word for saying what happens if.
\end{grammarlens}

\subsection*{Your turn}
\begin{itemize}
\raggedright
  \item \emph{Say it:} \emph{bahumānikkuka}
  \item \emph{Say it:} \emph{bahumānikkuka}, once more
  \item \emph{Say it:} \emph{kuṭṭikaḷ ā adhyāpakane vaḷare bahumānikkunnu.} — The children respect that teacher very much.
  \item \emph{Recall:} say the Malayalam for to confirm, then the Malayalam for to remind, then say \emph{bahumānikkuka} again
\end{itemize}

\begin{tcolorbox}[breakable,colback=teal!4,colframe=teal!35!black,title={Before you move on}]
What does \ml{ബഹുമാനിക്കുക} mean? (\textquotedblleft{}To look up to; to respect\textquotedblright{}.) Say it once more.
\end{tcolorbox}
\section[\texorpdfstring{\ml{തുല്യമായ} (tulyamāya)}{തുല്യമായ (tulyamāya)}]{\ml{തുല്യമായ} (tulyamāya) — equal}

\label{lesson:ML-C455-tulyamaya}



\begin{quote}
Before the new one: say the Malayalam for to remind, then the Malayalam for to look up to.
\end{quote}

\subsection*{You'll want to know: \ml{തുല്യമായ}}
\textbf{\ml{തുല്യമായ}} — \emph{tulyamāya} — \textquotedblleft{}equal\textquotedblright{}.

\begin{quote}
\ml{എല്ലാവർക്കും} \ml{തുല്യമായ} \ml{അവസരം} \ml{വേണം}.
\end{quote}

\begin{quote}
\emph{ellāvarkkuṁ tulyamāya avasaraṁ vēṇaṁ.}
\end{quote}

\begin{quote}
Everyone needs an equal chance.
\end{quote}

The adverb is \ml{തുല്യമായി}, equally.

\begin{cousinweb}
\textbf{\ml{തുല്യമായ}} (\emph{tulyamāya}) is built from Sanskrit \emph{tulya}, \textquotedblleft{}equal to, like\textquotedblright{} (CDIAL 5888), from the root \emph{tul}, \textquotedblleft{}to weigh\textquotedblright{}, with the adjective ending -\ml{ആയ}: of equal weight.
\end{cousinweb}

\begin{grammarlens}[title={What you've built}]
One more word for saying what happens if.
\end{grammarlens}

\subsection*{Your turn}
\begin{itemize}
\raggedright
  \item \emph{Say it:} \emph{tulyamāya}
  \item \emph{Say it:} \emph{tulyamāya}, once more
  \item \emph{Say it:} \emph{ellāvarkkuṁ tulyamāya avasaraṁ vēṇaṁ.} — Everyone needs an equal chance.
  \item \emph{Recall:} say the Malayalam for to remind, then the Malayalam for to look up to, then say \emph{tulyamāya} again
\end{itemize}

\begin{tcolorbox}[breakable,colback=teal!4,colframe=teal!35!black,title={Before you move on}]
What does \ml{തുല്യമായ} mean? (\textquotedblleft{}Equal\textquotedblright{}.) Say it once more.
\end{tcolorbox}
\section[\texorpdfstring{\ml{ധാർമ്മികമായ} (dhārmmikamāya)}{ധാർമ്മികമായ (dhārmmikamāya)}]{\ml{ധാർമ്മികമായ} (dhārmmikamāya) — ethical; moral}

\label{lesson:ML-C455-dharmmikamaya}



\begin{quote}
Before the new one: say the Malayalam for to look up to, then the Malayalam for equal.
\end{quote}

\subsection*{You'll want to know: \ml{ധാർമ്മികമായ}}
\textbf{\ml{ധാർമ്മികമായ}} — \emph{dhārmmikamāya} — \textquotedblleft{}ethical; moral\textquotedblright{}.

\begin{quote}
\ml{അത്} \ml{ധാർമ്മികമായ} \ml{ഒരു} \ml{തീരുമാനമായിരുന്നില്ല}.
\end{quote}

\begin{quote}
\emph{atŭ dhārmmikamāya oru tīrumānamāyirunnilla.}
\end{quote}

\begin{quote}
That was not a moral decision.
\end{quote}

The noun is \ml{ധാർമ്മികത}, morality.

\begin{cousinweb}
\textbf{\ml{ധാർമ്മികമായ}} (\emph{dhārmmikamāya}) is built from Sanskrit \emph{dhārmika}, \textquotedblleft{}righteous\textquotedblright{} (CDIAL 6798), made from \emph{dharma}, \textquotedblleft{}right conduct\textquotedblright{}, with the adjective ending -\ml{ആയ}.
\end{cousinweb}

\begin{grammarlens}[title={What you've built}]
One more word for saying what happens if.
\end{grammarlens}

\subsection*{Your turn}
\begin{itemize}
\raggedright
  \item \emph{Say it:} \emph{dhārmmikamāya}
  \item \emph{Say it:} \emph{dhārmmikamāya}, once more
  \item \emph{Say it:} \emph{atŭ dhārmmikamāya oru tīrumānamāyirunnilla.} — That was not a moral decision.
  \item \emph{Recall:} say the Malayalam for to look up to, then the Malayalam for equal, then say \emph{dhārmmikamāya} again
\end{itemize}

\begin{tcolorbox}[breakable,colback=teal!4,colframe=teal!35!black,title={Before you move on}]
What does \ml{ധാർമ്മികമായ} mean? (\textquotedblleft{}Ethical; moral\textquotedblright{}.) Say it once more.
\end{tcolorbox}
\section[\texorpdfstring{\ml{ഉപദ്രവിക്കുക} (upadravikkuka)}{ഉപദ്രവിക്കുക (upadravikkuka)}]{\ml{ഉപദ്രവിക്കുക} (upadravikkuka) — to harm; to harass}

\label{lesson:ML-C455-upadravikkuka}



\begin{quote}
Before the new one: say the Malayalam for equal, then the Malayalam for ethical.
\end{quote}

\subsection*{You'll want to know: \ml{ഉപദ്രവിക്കുക}}
\textbf{\ml{ഉപദ്രവിക്കുക}} — \emph{upadravikkuka} — \textquotedblleft{}to harm; to harass\textquotedblright{}.

\begin{quote}
\ml{ആ} \ml{നായ} \ml{ആരെയും} \ml{ഉപദ്രവിക്കില്ല}.
\end{quote}

\begin{quote}
\emph{ā nāya āreyuṁ upadravikkilla.}
\end{quote}

\begin{quote}
That dog will not harm anyone.
\end{quote}

Past: \ml{ഉപദ്രവിച്ചു}, \emph{upadraviccu}.

The noun is \ml{ഉപദ്രവം}, harm, trouble.

\begin{cousinweb}
\textbf{\ml{ഉപദ്രവിക്കുക}} (\emph{upadravikkuka}) is built from Sanskrit \emph{upadrava}, \textquotedblleft{}calamity, violence\textquotedblright{} (CDIAL 2180), with the Malayalam verb ending -\ml{ിക്കുക}.
\end{cousinweb}

\begin{grammarlens}[title={What you've built}]
One more word for saying what happens if.
\end{grammarlens}

\subsection*{Your turn}
\begin{itemize}
\raggedright
  \item \emph{Say it:} \emph{upadravikkuka}
  \item \emph{Say it:} \emph{upadravikkuka}, once more
  \item \emph{Say it:} \emph{ā nāya āreyuṁ upadravikkilla.} — That dog will not harm anyone.
  \item \emph{Recall:} say the Malayalam for equal, then the Malayalam for ethical, then say \emph{upadravikkuka} again
\end{itemize}

\begin{tcolorbox}[breakable,colback=teal!4,colframe=teal!35!black,title={Before you move on}]
What does \ml{ഉപദ്രവിക്കുക} mean? (\textquotedblleft{}To harm; to harass\textquotedblright{}.) Say it once more.
\end{tcolorbox}
